\documentclass[11pt]{article}

\usepackage[margin=1in]{geometry}
\usepackage{amsmath,amssymb,amsthm,mathtools}
\usepackage{bm}
\usepackage{booktabs}
\usepackage{longtable}
\usepackage{array}
\usepackage{hyperref}
\usepackage{graphicx}
\usepackage{enumitem}
\usepackage{xcolor}
\usepackage{ifthen}
\usepackage{caption}

\hypersetup{
  colorlinks=true,
  linkcolor=blue,
  citecolor=blue,
  urlcolor=blue,
  pdftitle={Verum Interior Gauge (VIG): Version 7 Final},
  pdfauthor={Stephen L. Nagy}
}

\newtheorem{theorem}{Theorem}[section]
\newtheorem{proposition}[theorem]{Proposition}
\newtheorem{corollary}[theorem]{Corollary}
\newtheorem{definition}[theorem]{Definition}
\newtheorem{remark}[theorem]{Remark}
\newtheorem{conjecture}[theorem]{Conjecture}

\newcommand{\Ccrit}{C_{\mathrm{crit}}}
\newcommand{\CR}{\mathrm{CR}}
\newcommand{\vig}{\mathrm{VIG}}
\newcommand{\dd}{\mathrm{d}}
\newcommand{\Hay}{\mathrm{Hay}}
\newcommand{\Bard}{\mathrm{Bar}}

\newcommand{\maybeincludegraphic}[3][]{%
  \IfFileExists{#2}{%
    \begin{figure}[htbp]
    \centering
    \includegraphics[#1]{#2}
    \caption{#3}
    \end{figure}
  }{}
}

\title{\Large\bfseries Verum Interior Gauge (VIG): Version 7 Final\\[4pt]
\large Comparative Critical-Growth Laws for Regular Black-Hole Interiors}
\author{Stephen L. Nagy}
\date{March 2026}

\begin{document}

\maketitle

\begin{abstract}
Version 7 Final presents the polished comparative form of the Verum Interior Gauge (\(\vig\)) program. The original motivating question was whether black-hole interiors admit a universal saturation clock of the form
\[
\tau_{\mathrm{sat}}\sim CMe^S.
\]
The geometry-first program led to a sharper and better-supported object: for each regular interior model studied so far, the near-critical maximal-slice family is governed by a model-dependent critical growth constant
\[
\Ccrit(M,u)=4\pi M^2\,\mathcal S(u),
\]
where \(u\) is a normalized deformation parameter. In this paper, the Hayward and Bardeen interiors are treated in a uniform formalism, each with an exact dimensionless law, a Schwarzschild anchor, a weak-deformation asymptotic, and a near-extremal asymptotic. The numerical family scans show the same structural law in both cases:
\[
\text{tail mean}\approx \Ccrit.
\]
The universality is therefore structural rather than exact-functional. Hayward and Bardeen belong to the same qualitative maximal-slice growth class, but they do not share the same deformation curve. In the full-edge overlap now available from the master driver, Bardeen remains below Hayward throughout and is approximately \(16.3\%\) lower in the normalized critical-growth law at the top of that overlap. The paper is tied explicitly to the source CSV products, the generating codes, and the master reproducibility driver.
\end{abstract}

\tableofcontents

\section{Introduction and motivation}

The starting point of this project was a broad physical intuition: if black-hole interiors are not merely passive continuations of exterior geometry but dynamical regions with their own large-scale structure, then one should be able to describe that structure with observables that scale cleanly across models. The original framing emphasized a possible \emph{Universal Clock}, inspired by the heuristic form
\[
\tau_{\mathrm{sat}}\sim CMe^S.
\]
That conjectural program remains interesting, but it is no longer the center of gravity of the work.

The geometry-first path produced a more immediate and more rigorously testable object. For spherically symmetric maximal slices in regular black-hole interiors, one finds a distinguished critical constant \(\Ccrit\) controlling the near-critical family. In the present framework, the operative law is not an entropy-weighted saturation statement, but a comparative maximal-slice growth law,
\[
V(v)\sim \Ccrit^{(\mathrm{model})}(M,u)\,v,
\]
where \(M\) is the black-hole mass and \(u\) is a dimensionless deformation parameter. The entire project now revolves around three linked goals:
\begin{enumerate}[leftmargin=2em]
\item derive exact expressions for \(\Ccrit\) in concrete regular metrics,
\item determine the weak-deformation and near-extremal asymptotics,
\item establish which aspects of the resulting growth law are universal and which are model-specific.
\end{enumerate}

The two regular models developed here are the Hayward metric and the Bardeen metric. The first is organized naturally by \(\lambda=\ell/M\), where \(\ell\) sets the Hayward core scale. The second is organized by \(\beta=g/M\), where \(g\) is the Bardeen magnetic-charge scale. Both reduce to Schwarzschild when the deformation vanishes. Both admit extremal endpoints. And both now have exact analytic \(\Ccrit\)-laws, numerical atlases, and verified asymptotics.

The vision of the \(\vig\) program is to turn regular black-hole interiors into a reproducible comparative cartography. In that spirit, this final version is tied directly to explicit code and exported data products: the Hayward unified framework, the Bardeen unified framework, the Bardeen dense edge scan, the full-edge Hayward--Bardeen comparison, and a master driver that aggregates the whole pipeline.

\section{Context and prior work}

The large-interior perspective for black holes was sharpened by Christodoulou and Rovelli, who showed that the maximal interior volume of a Schwarzschild black hole grows linearly with advanced time at the rate \(3\sqrt3\,\pi M^2\) \cite{ChristodoulouRovelli2015}. Bengtsson and Jakobsson clarified geometric aspects of these large interiors and emphasized how nontrivial the interior notion of volume becomes even in classical general relativity \cite{BengtssonJakobsson2015}. Subsequent work explored persistence of large interiors, thermodynamic questions, and the fate of such growth in modified scenarios \cite{Ong2015,Zhang2015}.

On the regular-black-hole side, Hayward proposed a nonsingular evaporating black-hole model that replaces the classical singularity by a finite core \cite{Hayward2006}. Bardeen's earlier regular black-hole construction remains historically foundational and has served as a template for many later nonsingular geometries \cite{Bardeen1968}. Modern discussions of nonsingular or geodesically complete black holes, including the status of regular cores and interior structure, may be found in \cite{Frolov2016,Carballo-Rubio2020}.

At the level of interior growth and late-time structure, exact results in lower-dimensional gravity also motivate caution about extrapolating from classical growth to true saturation. In particular, Iliesiu, Mezei, and S\'arosi showed that the volume of the black-hole interior in a solvable two-dimensional setting has a much richer late-time structure than a naive perpetual linear law would suggest \cite{IliesiuMezeiSarosi2022}. That kind of result keeps the original clock motivation alive, but it does not substitute for a fully four-dimensional geometric derivation.

The present work therefore stays disciplined: it does not claim a derivation of an entropy-weighted clock. Instead, it proves and verifies comparative geometric growth laws for regular interiors, using code and data products that can be rerun directly.

\section{Generic \texorpdfstring{$\vig$}{VIG} formalism}

\subsection{Metric and graph gauge}

Work in ingoing Eddington--Finkelstein coordinates,
\begin{equation}
 ds^2=-f(r)\,dv^2+2\,dv\,dr+r^2d\Omega^2.
\label{eq:efmetric}
\end{equation}
A spherically symmetric spacelike slice is written as a graph \(v=v(r)\). Its induced radial coefficient is
\begin{equation}
 \gamma_{rr}=-f(r)\,v'(r)^2+2v'(r),
\end{equation}
and the corresponding volume functional is
\begin{equation}
 V=\int 4\pi r^2\sqrt{-f(r)\,v'(r)^2+2v'(r)}\,dr.
\label{eq:volumefunctional}
\end{equation}

\subsection{Conserved quantity and reduced branch}

Because the reduced Lagrangian does not depend explicitly on \(v\), there is a conserved quantity,
\begin{equation}
 C=\frac{4\pi r^2\bigl(1-f(r)v'(r)\bigr)}{\sqrt{-f(r)\,v'(r)^2+2v'(r)}}.
\label{eq:genericC}
\end{equation}
Defining
\begin{equation}
 Q(r;C)=C^2+(4\pi r^2)^2f(r),
\label{eq:genericQ}
\end{equation}
the smooth maximal-slice branch takes the compact form
\begin{equation}
 v'(r)=\frac{1-C/\sqrt{Q(r;C)}}{f(r)},
\label{eq:genericvprime}
\end{equation}
while the volume integrand becomes
\begin{equation}
 \frac{dV}{dr}=\frac{(4\pi r^2)^2}{\sqrt{Q(r;C)}}.
\label{eq:genericdVdr}
\end{equation}

This reduction is the technical backbone of the whole program. Every model-specific result in the paper is obtained by specifying \(f(r)\), determining the admissible branch, and locating the maximizing point of the induced critical quantity.

\subsection{Critical constant}

The admissibility condition is \(Q(r;C)\ge 0\) on the physical interval \(0<r<r_+\). The critical constant is therefore
\begin{equation}
\boxed{
\Ccrit^2(M,\text{deformation})=
\max_{0<r<r_+}\Bigl[-(4\pi r^2)^2 f(r)\Bigr].
}
\label{eq:genericCcrit}
\end{equation}
This quantity is the central invariant of the \(\vig\) framework.

\section{Schwarzschild benchmark}

For Schwarzschild,
\[
f_{\mathrm{Schw}}(r)=1-\frac{2M}{r}.
\]
The critical point is
\[
r_*=\frac{3M}{2},
\]
and the critical constant is exactly
\begin{equation}
\boxed{
\Ccrit(M,0)=3\sqrt3\,\pi M^2.
}
\label{eq:SchwarzschildCcrit}
\end{equation}
This is precisely the Christodoulou--Rovelli growth coefficient \cite{ChristodoulouRovelli2015}. It serves as the anchor of the comparative theory: every regular model studied here reduces continuously to this value at zero deformation.

\section{Hayward model}

\subsection{Metric and exact reduction}

For Hayward,
\begin{equation}
 f_{\Hay}(r)=1-\frac{2Mr^2}{r^3+2M\ell^2}.
\label{eq:HaywardMetric}
\end{equation}
Introduce the dimensionless variables
\[
x=\frac{r}{M},\qquad \lambda=\frac{\ell}{M}.
\]
Then
\begin{equation}
 f_{\Hay}(x,\lambda)=1-\frac{2x^2}{x^3+2\lambda^2}.
\end{equation}
Defining
\begin{equation}
\Phi_\lambda(x):=-x^4 f_{\Hay}(x,\lambda)=
\frac{x^4\left(2x^2-x^3-2\lambda^2\right)}{x^3+2\lambda^2},
\end{equation}
the exact Hayward critical law is
\begin{equation}
\boxed{
\Ccrit^{\Hay}(M,\ell)=4\pi M^2\,\mathcal C_{\Hay}(\lambda),
\qquad
\mathcal C_{\Hay}(\lambda)=\sqrt{\max_{0<x<x_+}\Phi_\lambda(x)}.
}
\label{eq:HaywardExactLaw}
\end{equation}

\subsection{Critical and horizon equations}

Let \(x_*(\lambda)\) maximize \(\Phi_\lambda(x)\). Then the critical point satisfies
\begin{equation}
\boxed{
2x^6-3x^5+8\lambda^2x^3-12\lambda^2x^2+8\lambda^4=0.
}
\label{eq:HaywardCriticalPoly}
\end{equation}
The outer horizon is determined by
\begin{equation}
\boxed{
x_+^3-2x_+^2+2\lambda^2=0.
}
\label{eq:HaywardHorizonCubic}
\end{equation}
At extremality,
\begin{equation}
\boxed{
\lambda_c=\frac{4}{3\sqrt3},
\qquad
x_c=\frac43.
}
\label{eq:HaywardExtremal}
\end{equation}

\subsection{Hayward asymptotics}

Near \(\lambda=0\),
\begin{equation}
 x_*(\lambda)=\frac32-\frac{128}{243}\lambda^4+\mathcal O(\lambda^6),
\end{equation}
and
\begin{equation}
\boxed{
\Ccrit^{\Hay}(M,\ell)=3\sqrt3\,\pi M^2
\left[
1-\frac{32}{27}\left(\frac{\ell}{M}\right)^2
+\mathcal O\!\left(\frac{\ell^4}{M^4}\right)
\right].
}
\label{eq:HaywardSmallLambda}
\end{equation}

Near extremality, with
\[
t=\lambda_c^2-\lambda^2,
\]
one finds
\begin{align}
 x_+(\lambda)
 &=\frac43+\sqrt{t}-\frac14 t+\frac{5}{32}t^{3/2}+\mathcal O(t^2),\\
 x_*(\lambda)
 &=\frac43+\frac34 t-\frac{135}{64}t^2+\mathcal O(t^3),
\end{align}
and the critical-growth law becomes
\begin{equation}
\boxed{
\Ccrit^{\Hay}(M,\ell)=
\frac{16\pi}{3}M^2\sqrt{t}
+3\pi M^2 t^{3/2}
-\frac{99\pi}{32}M^2 t^{5/2}
+\mathcal O(t^{7/2}).
}
\label{eq:HaywardNearExt}
\end{equation}

\subsection{Hayward numerical status}

The exact Hayward analytic law was checked directly against the numerical scan products exported by the Hayward unified framework. The source products \texttt{vig\_interior\_atlas.csv}, \texttt{vig\_interior\_atlas\_unified.csv}, and \texttt{vig\_unified\_master\_summary.csv} supply the atlas, summary, and asymptotic-validation layers. In the master results, Hayward's tail-to-critical ratio averages approximately
\[
\frac{\text{tail mean}}{\Ccrit^{\Hay}}\approx 1.001206,
\]
showing that the tail law is not merely qualitative but numerically tight.

\section{Bardeen model}

\subsection{Metric and exact reduction}

For Bardeen,
\begin{equation}
 f_{\Bard}(r)=1-\frac{2Mr^2}{(r^2+g^2)^{3/2}}.
\label{eq:BardeenMetric}
\end{equation}
Introduce
\[
x=\frac{r}{M},\qquad \beta=\frac{g}{M}.
\]
Then
\begin{equation}
 f_{\Bard}(x,\beta)=1-\frac{2x^2}{(x^2+\beta^2)^{3/2}}.
\end{equation}
Defining
\begin{equation}
\Phi_\beta(x):=-x^4 f_{\Bard}(x,\beta)=
\frac{x^4\bigl(2x^2-(x^2+\beta^2)^{3/2}\bigr)}{(x^2+\beta^2)^{3/2}},
\end{equation}
the exact Bardeen critical law is
\begin{equation}
\boxed{
\Ccrit^{\Bard}(M,g)=4\pi M^2\,\mathcal C_{\Bard}(\beta),
\qquad
\mathcal C_{\Bard}(\beta)=\sqrt{\max_{0<x<x_+}\Phi_\beta(x)}.
}
\label{eq:BardeenExactLaw}
\end{equation}

\subsection{Auxiliary variable and exact equations}

For Bardeen, the reduction becomes especially transparent in the auxiliary variable
\[
y=\sqrt{x^2+\beta^2}.
\]
Then \(x^2=y^2-\beta^2\), and the critical-point equation collapses to
\begin{equation}
\boxed{
2y^5-3y^4+3\beta^4=0.
}
\label{eq:BardeenCriticalY}
\end{equation}
The horizon equation becomes
\begin{equation}
\boxed{
y^3-2y^2+2\beta^2=0.
}
\label{eq:BardeenHorizonY}
\end{equation}
Thus the exact Bardeen law is obtained by solving \eqref{eq:BardeenCriticalY} for the physical root \(y_*(\beta)\), reconstructing \(x_*(\beta)=\sqrt{y_*(\beta)^2-\beta^2}\), and evaluating
\[
\Ccrit^{\Bard}(M,g)=4\pi M^2\,x_*(\beta)^2\sqrt{-f_{\Bard}(x_*(\beta),\beta)}.
\]

\subsection{Bardeen extremality and asymptotics}

The extremal threshold is
\begin{equation}
\boxed{
\beta_c=\frac{4}{3\sqrt3}.
}
\label{eq:BardeenBetac}
\end{equation}
At the extremal horizon,
\[
y_c=\frac43,
\qquad
x_{+,c}^2=\frac{32}{27},
\qquad
x_{+,c}=\frac{4\sqrt6}{9}.
\]

Near \(\beta=0\),
\begin{equation}
 y_*(\beta)=\frac32-\frac{8}{27}\beta^4+\mathcal O(\beta^8),
\end{equation}
and
\begin{equation}
\boxed{
\Ccrit^{\Bard}(M,g)=3\sqrt3\,\pi M^2
\left[
1-\frac43\left(\frac{g}{M}\right)^2
+\mathcal O\!\left(\frac{g^4}{M^4}\right)
\right].
}
\label{eq:BardeenSmallBeta}
\end{equation}

Near extremality, with
\[
t=\beta_c^2-\beta^2,
\]
one finds
\begin{align}
 y_+(\beta)
 &=\frac43+\sqrt{t}-\frac14 t+\frac{5}{32}t^{3/2}+\mathcal O(t^2),\\
 y_*(\beta)
 &=\frac43+\frac98 t-\frac{3645}{512}t^2+\mathcal O(t^3),
\end{align}
and the leading critical-growth law becomes
\begin{equation}
\boxed{
\Ccrit^{\Bard}(M,g)=
\frac{16\sqrt6\,\pi}{9}M^2\sqrt{t}
+\frac{21\sqrt6\,\pi}{8}M^2 t^{3/2}
-\frac{13689\sqrt6\,\pi}{2048}M^2 t^{5/2}
+\mathcal O(t^{7/2}).
}
\label{eq:BardeenNearExt}
\end{equation}

\subsection{Bardeen numerical status}

The exact Bardeen law is verified by the source products \texttt{vig\_bardeen\_atlas.csv}, \texttt{vig\_bardeen\_master\_summary.csv}, \texttt{vig\_bardeen\_edge\_atlas.csv}, and \texttt{vig\_bardeen\_edge\_master\_summary.csv}. The dense edge scan shows that the near-extremal expansion enters its expected regime near the edge and that the tail law remains numerically tight throughout. In the master results, the mean Bardeen tail-to-critical ratio is approximately
\[
\frac{\text{tail mean}}{\Ccrit^{\Bard}}\approx 1.001258.
\]
The Hayward and Bardeen tail laws are therefore comparably stable.

\section{Comparative full-edge results}

\subsection{Universal structure}

The following pattern now appears in every model studied:

\begin{proposition}[Comparative \(\vig\) universality, empirical form]
For Schwarzschild, Hayward, and Bardeen, the near-critical maximal-slice family obeys
\[
V(v)\sim \Ccrit^{(\mathrm{model})}(M,u)\,v,
\]
with the numerical tail satisfying
\[
\text{tail mean}\approx \Ccrit^{(\mathrm{model})}.
\]
The Schwarzschild anchor is recovered at vanishing deformation:
\[
\Ccrit^{(\mathrm{model})}(M,0)=3\sqrt3\,\pi M^2.
\]
\end{proposition}

The common qualitative geometry is equally clear:
\begin{itemize}[leftmargin=2em]
\item the outer horizon shrinks as deformation grows,
\item the critical slice radius also shrinks,
\item the admissible maximal-slice channel compresses,
\item the normalized critical-growth law decreases monotonically.
\end{itemize}

\subsection{Model-specific structure}

The universality does \emph{not} mean identical deformation curves. At equal normalized deformation \(u\),
\begin{equation}
\frac{\Ccrit^{\Bard}}{3\sqrt3\,\pi M^2}<
\frac{\Ccrit^{\Hay}}{3\sqrt3\,\pi M^2}
\qquad (u>0)
\label{eq:BardeenBelowHayward}
\end{equation}
through the full overlapping comparison range available in the master outputs.

The decisive new result is the full-edge comparison. Using the aggregated products \texttt{master\_full\_edge\_common\_grid.csv} and \texttt{master\_full\_edge\_combined\_atlas.csv}, the comparison now extends to the top of the overlapping normalized interval. The resulting behavior is unambiguous:
\begin{itemize}[leftmargin=2em]
\item At \(u=0\), the two models coincide, as required by the Schwarzschild anchor.
\item As \(u\) increases, Bardeen remains below Hayward at every sampled point in the normalized critical-growth law.
\item The difference is mild at small deformation, but accelerates strongly near the edge.
\item By the top of the overlap, the ratio has fallen to approximately
\[
\frac{\Ccrit^{\Bard}/(3\sqrt3\pi M^2)}{\Ccrit^{\Hay}/(3\sqrt3\pi M^2)}\approx 0.837,
\]
so Bardeen's normalized critical-growth law is about \(16.3\%\) lower than Hayward's there.
\end{itemize}

This goes well beyond the earlier partial-overlap comparison. The two models remain in the same qualitative class, but they do not remain quantitatively close near the edge.

\begin{remark}
The weak-deformation asymptotics already hint at the correct hierarchy:
\[
\Ccrit^{\Hay}=3\sqrt3\,\pi M^2\left[1-\frac{32}{27}u^2+\cdots\right],
\qquad
\Ccrit^{\Bard}=3\sqrt3\,\pi M^2\left[1-\frac43u^2+\cdots\right].
\]
Since \(\tfrac43>\tfrac{32}{27}\), Bardeen bends downward faster near \(u=0\). The full-edge comparison shows that this early tendency amplifies rather than washes out.
\end{remark}

\subsection{Geometric comparison}

The geometric comparison is just as important as the growth-law comparison.
\begin{itemize}[leftmargin=2em]
\item Bardeen has smaller \(x_+\) than Hayward at equal normalized deformation.
\item Bardeen has smaller \(x_*\) than Hayward at equal normalized deformation.
\item The entire Bardeen structure is therefore pushed inward more strongly.
\end{itemize}

The most interesting twist lies in the interior gap
\[
x_+-x_*.
\]
Although both \(x_+\) and \(x_*\) are smaller in Bardeen, the comparative plots show that the Bardeen gap is slightly \emph{larger} than the Hayward gap over the overlap region. So Bardeen is not merely applying a uniform rescaling of the Hayward geometry. It is pushing the whole structure inward more strongly while preserving a slightly wider separation between the outer horizon and the critical slice.

This makes the comparative result geometrically vivid: the two models share the same type of interior growth law, but the way they deform the interior channel is not the same.

\section{Reproducibility architecture}

Version 7 Final is intended to be reproducible from a finite and explicit code--data chain.

\subsection{Source codes}

The core generating codes are:
\begin{center}
\begin{tabular}{>{\ttfamily}p{0.37\linewidth} p{0.53\linewidth}}
\toprule
File & Role \\
\midrule
vig\_unified\_framework\_v4.py & Hayward unified framework; generates Hayward atlas and Hayward master summaries. \\
vig\_bardeen\_unified\_v1.py & Bardeen unified framework; generates Bardeen atlas and summary products. \\
vig\_bardeen\_edge\_scan\_v1.py & Dense Bardeen edge scan and asymptotic validation near extremality. \\
vig\_hayward\_bardeen\_full\_edge\_comparison.py & Full-edge comparative overlay between Hayward and Bardeen. \\
vig\_master\_driver.py & Top-level aggregation layer; exports master CSVs and final figure suite. \\
universal\_clock.py & Original clock bookkeeping and invariant summary. \\
\bottomrule
\end{tabular}
\end{center}

\subsection{Source data products}

The principal source CSVs are:
\begin{center}
\begin{tabular}{>{\ttfamily}p{0.43\linewidth} p{0.47\linewidth}}
\toprule
CSV product & Role \\
\midrule
vig\_analytic\_ccrit\_validation.csv & Hayward analytic/numeric validation layer. \\
vig\_interior\_atlas.csv & Hayward atlas. \\
vig\_interior\_atlas\_unified.csv & Hayward unified atlas export. \\
vig\_unified\_master\_summary.csv & Hayward unified master summary. \\
vig\_bardeen\_atlas.csv & Bardeen atlas. \\
vig\_bardeen\_master\_summary.csv & Bardeen summary layer. \\
vig\_bardeen\_edge\_atlas.csv & Dense Bardeen edge atlas. \\
vig\_bardeen\_edge\_master\_summary.csv & Dense Bardeen edge summary layer. \\
vig\_hayward\_bardeen\_full\_edge\_comparison.csv & Combined comparison export. \\
vig\_hayward\_bardeen\_full\_edge\_common\_grid.csv & Common interpolation grid for full-edge comparison. \\
\bottomrule
\end{tabular}
\end{center}

\subsection{Master outputs}

The master driver compiles these layers into the final aggregation products:
\begin{center}
\begin{tabular}{>{\ttfamily}p{0.43\linewidth} p{0.47\linewidth}}
\toprule
Master output & Role \\
\midrule
master\_universal\_clock\_summary.csv & Clock bookkeeping summary. \\
master\_hayward\_atlas.csv & Final Hayward atlas used for cross-model comparison. \\
master\_bardeen\_atlas.csv & Final Bardeen atlas used for cross-model comparison. \\
master\_full\_edge\_common\_grid.csv & Canonical common grid for the full-edge overlap. \\
master\_full\_edge\_combined\_atlas.csv & Combined model atlas for comparison. \\
master\_multimodel\_summary.csv & Compact summary of model ranges and tail-tracking statistics. \\
master\_manifest.csv & Artifact manifest for the aggregated run. \\
\bottomrule
\end{tabular}
\end{center}

The manuscript figures are tied directly to the corresponding exported PNGs from this pipeline. In other words, the final paper is not merely inspired by the code; it is derived from an explicit code path and a finite set of named outputs.

\section{Selected numerical outputs}

\maybeincludegraphic[width=0.82\linewidth]{master_compare_edge_Ccrit_over_CR_overlay.png}{Hayward vs. Bardeen overlay of the normalized critical-growth law over the full overlapping normalized-deformation range.}

\maybeincludegraphic[width=0.82\linewidth]{master_compare_edge_ratio_bardeen_over_hayward.png}{Bardeen-to-Hayward ratio of the normalized critical-growth law over the full-edge overlap. The ratio falls from the Schwarzschild anchor to approximately \(0.837\) near the top of the overlap.}

\maybeincludegraphic[width=0.82\linewidth]{master_compare_edge_shape_function_overlay.png}{Comparison of the shape function \(\Ccrit/(4\pi M^2)\) for Hayward and Bardeen. The separation is not a normalization artifact; it is present directly in the underlying dimensionless law.}

\maybeincludegraphic[width=0.82\linewidth]{master_compare_edge_xplus_overlay.png}{Outer-horizon comparison between Hayward and Bardeen over the full-edge overlap. Bardeen contracts the outer horizon more strongly.}

\maybeincludegraphic[width=0.82\linewidth]{master_compare_edge_xstar_overlay.png}{Critical-slice comparison between Hayward and Bardeen over the full-edge overlap. Bardeen pushes the critical slice inward more strongly.}

\maybeincludegraphic[width=0.82\linewidth]{master_compare_edge_gap_overlay.png}{Comparison of the interior channel gap \(x_+-x_*\). Although Bardeen has smaller \(x_+\) and smaller \(x_*\), its gap is slightly larger over the overlap.}

\maybeincludegraphic[width=0.82\linewidth]{vig_bardeen_exact_vs_asymptotics.png}{Bardeen exact law versus weak-deformation and near-extremal asymptotics in the dense edge scan.}

\maybeincludegraphic[width=0.82\linewidth]{vig_bardeen_ccrit_vs_tail.png}{Bardeen edge scan: the near-critical family tail continues to track \(\Ccrit\) toward extremality.}

\section{Discussion and scope}

The central accomplishment of the present program is not the derivation of a universal saturation time. It is the discovery and verification of a comparative geometric law for regular interiors. The strongest validated statement is
\[
V(v)\sim \Ccrit^{(\mathrm{model})}(M,u)\,v,
\qquad
\text{tail mean}\approx \Ccrit^{(\mathrm{model})}.
\]
That statement is supported by exact reduction, asymptotic analysis, family scans, cross-model comparison, and a master reproducibility pipeline.

The original clock program remains conceptually alive but formally subordinate. Any future argument for a genuine saturation clock must now explain how it emerges from, modifies, or ultimately replaces the validated geometric critical-growth law. In that sense, the clock conjecture has become a \emph{downstream} problem rather than the primary theorem candidate.

The comparative outcome is also physically suggestive. Hayward and Bardeen both regularize the black-hole interior, but they do so differently enough that the resulting interior cartographies visibly diverge near the edge. The shared class is real; the model-specific deformation is also real. This is exactly the kind of result a comparative program should produce.

\section{Conclusion}

Version 7 Final is the polished culmination of the current \(\vig\) program.

The main result is no longer a one-model speculation about a possible interior clock. It is a two-model universality statement about near-critical maximal-slice growth:
\[
\boxed{
V(v)\sim \Ccrit^{(\mathrm{model})}(M,u)\,v,
\qquad
\text{tail mean}\approx \Ccrit^{(\mathrm{model})}.
}
\]

What is universal:
\begin{itemize}[leftmargin=2em]
\item Schwarzschild anchoring at zero deformation,
\item monotone suppression of the critical-growth law,
\item compression of the admissible interior channel,
\item numerical tracking of the tail by the model-specific \(\Ccrit\).
\end{itemize}

What is model-specific:
\begin{itemize}[leftmargin=2em]
\item the exact deformation curve \(\mathcal C(u)\),
\item the detailed horizon and critical-slice radii,
\item the strength of suppression at fixed normalized deformation,
\item the edge behavior of the comparative ratio.
\end{itemize}

Hayward and Bardeen therefore belong to the same qualitative interior-growth class, while Bardeen deforms that class more strongly. The full-edge comparison shows that this is not a small perturbative effect confined to weak deformation; it becomes substantial near the overlap boundary and is visible simultaneously in the growth law, the shape function, and the horizon--slice geometry. This is the strongest, cleanest, and most reproducible statement the project has produced so far.

\appendix

\section{Collected exact laws}

\subsection{Hayward}
\[
\Ccrit^{\Hay}(M,\ell)=4\pi M^2\,\mathcal C_{\Hay}(\lambda),
\qquad
\lambda=\frac{\ell}{M}.
\]
\[
2x^6-3x^5+8\lambda^2x^3-12\lambda^2x^2+8\lambda^4=0,
\qquad
x_+^3-2x_+^2+2\lambda^2=0.
\]

\subsection{Bardeen}
\[
\Ccrit^{\Bard}(M,g)=4\pi M^2\,\mathcal C_{\Bard}(\beta),
\qquad
\beta=\frac{g}{M}.
\]
\[
2y^5-3y^4+3\beta^4=0,
\qquad
y^3-2y^2+2\beta^2=0.
\]

\section{Collected weak-deformation laws}

\[
\Ccrit^{\Hay}(M,\ell)=3\sqrt3\,\pi M^2
\left[
1-\frac{32}{27}\left(\frac{\ell}{M}\right)^2+\cdots
\right]
\]
\[
\Ccrit^{\Bard}(M,g)=3\sqrt3\,\pi M^2
\left[
1-\frac43\left(\frac{g}{M}\right)^2+\cdots
\right]
\]

\section{Collected near-extremal laws}

For Hayward,
\[
\Ccrit^{\Hay}(M,\ell)=
\frac{16\pi}{3}M^2\sqrt{t}
+3\pi M^2 t^{3/2}
-\frac{99\pi}{32}M^2 t^{5/2}+\cdots,
\qquad
t=\lambda_c^2-\lambda^2.
\]

For Bardeen,
\[
\Ccrit^{\Bard}(M,g)=
\frac{16\sqrt6\,\pi}{9}M^2\sqrt{t}
+\frac{21\sqrt6\,\pi}{8}M^2 t^{3/2}
-\frac{13689\sqrt6\,\pi}{2048}M^2 t^{5/2}+\cdots,
\qquad
t=\beta_c^2-\beta^2.
\]

\begin{thebibliography}{99}

\bibitem{ChristodoulouRovelli2015}
M.~Christodoulou and C.~Rovelli,
``How big is a black hole?'',
\emph{Phys. Rev. D} \textbf{91}, 064046 (2015),
arXiv:1411.2854 [gr-qc].

\bibitem{BengtssonJakobsson2015}
I.~Bengtsson and E.~Jakobsson,
``Black holes: their large interiors,'',
\emph{Mod. Phys. Lett. A} \textbf{30}, 1550103 (2015),
arXiv:1502.01907 [gr-qc].

\bibitem{Hayward2006}
S.~A.~Hayward,
``Formation and evaporation of nonsingular black holes,'',
\emph{Phys. Rev. Lett.} \textbf{96}, 031103 (2006),
gr-qc/0506126.

\bibitem{Bardeen1968}
J.~M.~Bardeen,
``Non-singular general-relativistic gravitational collapse,'',
in \emph{Proceedings of GR5} (Tbilisi, USSR, 1968), p.~174.

\bibitem{Frolov2016}
V.~P.~Frolov,
``Notes on nonsingular models of black holes,'',
\emph{Phys. Rev. D} \textbf{94}, 104056 (2016),
arXiv:1609.01758 [gr-qc].

\bibitem{Carballo-Rubio2020}
R.~Carballo-Rubio, F.~Di~Filippo, S.~Liberati, and M.~Visser,
``Geodesically complete black holes,'',
\emph{Phys. Rev. D} \textbf{101}, 084047 (2020),
arXiv:1911.11200 [gr-qc].

\bibitem{Ong2015}
Y.~C.~Ong,
``The persistence of the large volumes in black holes,'',
\emph{Gen. Relativ. Gravit.} \textbf{47}, 88 (2015),
arXiv:1503.08245 [gr-qc].

\bibitem{Zhang2015}
Y.~Zhang,
``Entropy in the interior of a black hole and thermodynamics,'',
\emph{Gen. Relativ. Gravit.} \textbf{47}, 14 (2015),
arXiv:1407.6331 [gr-qc].

\bibitem{IliesiuMezeiSarosi2022}
L.~V.~Iliesiu, M.~Mezei, and G.~S\'arosi,
``The volume of the black hole interior at late times,'',
\emph{JHEP} \textbf{07}, 073 (2022),
arXiv:2107.06286 [hep-th].

\bibitem{HawkingEllis1973}
S.~W.~Hawking and G.~F.~R.~Ellis,
\emph{The Large Scale Structure of Space-Time}
(Cambridge University Press, 1973).

\bibitem{Hawking1975}
S.~W.~Hawking,
``Particle creation by black holes,'',
\emph{Commun. Math. Phys.} \textbf{43}, 199 (1975).

\bibitem{BardeenPressTeukolsky1972}
J.~M.~Bardeen, W.~H.~Press, and S.~A.~Teukolsky,
``Rotating black holes: locally nonrotating frames, energy extraction, and scalar synchrotron radiation,'',
\emph{Astrophys. J.} \textbf{178}, 347 (1972).

\bibitem{Nagy2025BHRuler}
S.~L.~Nagy,
\emph{BHRuler: A Scale-Invariant Ruler for Black Holes},
GitHub/Zenodo (2025--2026),
\url{https://github.com/SNagy3/BHRuler},
doi:10.5281/zenodo.18894798.

\end{thebibliography}

\end{document}
